\documentclass[11pt]{article}
\usepackage[a4paper,margin=18mm]{geometry}
\usepackage{fontspec}
\setmainfont{Latin Modern Roman}
\usepackage{amsmath,amssymb,amsthm,amscd,graphicx,color}
\usepackage{xurl}
\usepackage[unicode,hidelinks]{hyperref}
\allowdisplaybreaks
\setlength\emergencystretch{3em}

\makeatletter
\newtheorem{theorem}{Theorem}
\newtheorem*{theorem*}{Theorem}
\newtheorem{lemma}{Lemma}
\newtheorem*{lemma*}{Lemma}
\newtheorem{corollary}{Corollary}
\newtheorem*{corollary*}{Corollary}
\newtheorem{proposition}{Proposition}
\newtheorem*{proposition*}{Proposition}
\newtheorem{conjecture}{Conjecture}
\newtheorem*{conjecture*}{Conjecture}
{\theoremstyle{definition} \newtheorem{definition}{Definition}
\newtheorem*{definition*}{Definition}
\newtheorem{example}{Example}
\newtheorem*{example*}{Example}
\newtheorem{remark}{Remark}
\newtheorem*{remark*}{Remark}
\newtheorem{note}{Note}
\newtheorem*{note*}{Note}
}

\makeatother
\newtheorem{Theorem}{Theorem}[section]
\newtheorem{Corollary}[Theorem]{Corollary}
\newtheorem{Lemma}[Theorem]{Lemma}
\newtheorem{Proposition}[Theorem]{Proposition}
{\theoremstyle{definition}
\newtheorem{Definition}[Theorem]{Definition}
\newtheorem{Note}[Theorem]{Note}
\newtheorem{Example}[Theorem]{Example}
\newtheorem{Remark}[Theorem]{Remark}
}

\begin{document}
\begin{center}\Large Equivalence of the normalized four-term BKP addition identity with the formal BKP residue hierarchy on a domain where tau and the displayed shifted denominators are invertible\end{center}
\noindent\textbf{Stable ID:} 2031\par
\noindent\textbf{Authors:} Yoko Shigyo\par
\noindent\textbf{Original work:} On Addition Formulae of KP, mKP and BKP Hierarchies\par
\noindent\textbf{Version:} Official SIGMA 9 (2013) article 035, DOI 10.3842/SIGMA.2013.035; journal PDF and native TeX acquired 2026-10-01, exact bytes pinned by SHA256\par
\noindent\textbf{Selected task scope:} Full Theorem3 in its normalized rational/formal setting: identity32 for arbitrary admissible Miwa parameters is equivalent to residue hierarchy28 in odd times. Tau and each displayed shifted tau/Miwa denominator are taken invertible wherever ratios are used. The residue is the explicit coefficient-extraction operator, not an analytic contour theorem; no extension across tau zeros is asserted.\par
\noindent\textbf{Difficulty:} H2: The proof builds all odd and even multipoint tau ratios from the four-term relation through an inductive Pfaffian construction and a Pfaffian Sylvester identity. Pfaffian Plucker relations recover the complete addition system, and the weighted power-sum specialization argument then recovers every coefficient of the formal residue hierarchy.\par
\noindent\textbf{Editorial boundary note (not author text):} The complete original Theorem 3 BKP equivalence is retained with the exact normalized four-term identity, formal residue hierarchy and all required author proof. The weighted-specialization Proposition 1 proof, odd Miwa times, odd/even addition identities, shifted four-point base, full Pfaffian induction and explicit Appendix C Plucker formulas are supplied. Sylvester Pfaffian identity is genuinely prior standard background and recorded in Appendix B. The domain is formal/localized where tau and every displayed shifted denominator and Miwa parameter factor are invertible; no global zero-locus or analytic-contour claim is added. The author-described even-n repetition and direct substitution are bounded uses of the supplied formulas. Original journal lowercase independent theorem/proposition/example counters are restored at the BKP boundary. KP/mKP restatements are not selected as additional units.\par
\noindent\textbf{Source:} \url{https://sigma-journal.com/2013/035/}\quad DOI: \url{https://doi.org/10.3842/SIGMA.2013.035}\par
\noindent\textbf{License and attribution:} Copyright retained by the named authors. Published by SIGMA under CC BY 4.0: \url{https://creativecommons.org/licenses/by/4.0/}. Source: \url{https://sigma-journal.com/about.html}. Adaptation consists of standalone excerpt formatting, cover and numbering restoration; original author language and mathematical text are retained.\par
\medskip\hrule\medskip\noindent\textbf{Author text begins below.}\par
\setcounter{section}{1}
\setcounter{subsection}{0}
\setcounter{subsubsection}{0}
\setcounter{Theorem}{0}
\setcounter{equation}{4}
% Original source lines 185--316
\section[The addition formula for the $\tau$-function of the KP hierarchy]{The addition formula for the $\boldsymbol{\tau}$-function of the KP hierarchy} \label{section2}

Let
\begin{gather*}
[\alpha]=\left(\alpha,\frac{\alpha^2}{2},\frac{\alpha^3}{3},\dots\right),
\qquad
\xi(t,\lambda)=\sum_{n=1}^{\infty}t_n\lambda^n,
\qquad
t=(t_1,t_2,t_3,\dots).
\end{gather*}
The KP hierarchy is a~system of equations for a~function $\tau(t)$ \cite{DJKM1,JM1, MJD1} given by
\begin{gather}
\oint e^{\xi(t'-t,\lambda)}\tau\big(t'-\big[\lambda^{-1}\big]\big)\tau\big(t+\big[\lambda^{-1}\big]\big)\frac{d\lambda}{2\pi i}=0.
\label{eq:1}
\end{gather}
Here $\oint$ means a~formal algebraic operator extracting the coef\/f\/icient of $z^{-1}$ of Laurent series:
\begin{gather*}%\label{eq:2}
\oint\frac{dz}{2\pi i}\sum_{n=-\infty}^{\infty}a_n z^n=a_{-1}.
\end{gather*}
Set $t=x+y$, $t'=x-y$.
Then~\eqref{eq:1} becomes
\begin{gather}
\oint e^{-2\xi(y,\lambda)}\tau\big(x-y-\big[\lambda^{-1}\big]\big)\tau\big(x+y+\big[\lambda^{-1}\big]\big)\frac{d\lambda}{2\pi i}=0.
\label{eq:3}
\end{gather}
For an integer $m\ge2$, set
\begin{gather}
y=\frac{1}{2}\left(\sum_{i=1}^{m-1}[\beta_i]-\sum_{i=1}^{m+1}[\alpha_i]\right)
\label{eq:22}
\end{gather}
in~\eqref{eq:3}. Then it becomes
\begin{gather}
\oint\exp\left({-\xi\left(\displaystyle\sum_{i=1}^{m-1}[\beta_i]-\sum_{i=1}^{m+1}[\alpha_i],\lambda\right)}\right)
\tau\left(x-\frac{1}{2}\left(\sum_{i=1}^{m-1}[\beta_i]-\sum_{i=1}^{m+1}[\alpha_i]\right)-\big[\lambda^{-1}
\big]\right)\nonumber
\\
\qquad
{}\times\tau\left(x+\frac{1}{2}\left(\sum_{i=1}^{m-1}[\beta_i]-\sum_{i=1}^{m+1}[\alpha_i]\right)+\big[\lambda^{-1}
\big]\right)\frac{d\lambda}{2\pi i}=0.
\label{eq:4}
\end{gather}
By virtue of the identity
\begin{gather*}
\sum_{n=1}^{\infty}\frac{x^n}{n}=-\log(1-x),
\end{gather*}
the exponential factor in~\eqref{eq:4} reduces to a~rational function of $\lambda$, $\alpha_i$, $\beta_i$ as
\begin{gather*}
\exp\left(-\xi\left(\sum_{i=1}^{m-1}[\beta_i]-\sum_{i=1}^{m+1}[\alpha_i],\lambda\right)\right)
=\frac{\prod\limits_{i=1}^{m-1}(1-\beta_i\lambda)}{\prod\limits_{i=1}^{m+1}
(1-\alpha_i\lambda)}.
\end{gather*}
Computing the integral by taking residues at $\lambda=\alpha_i^{-1}$, $1\le i\le m+1$ \cite{M1} and
shifting the variable $x$ as
\begin{gather*}
x\rightarrow x+\frac{1}{2}\left(\sum_{i=1}^{m-1}[\beta_i]+\sum_{i=1}^{m+1}[\alpha_i]\right),
\end{gather*}
we get the following addition formulae of the $\tau$-function~\cite{SS1}
\begin{gather}
\sum_{i=1}^{m+1}(-1)^{i-1}\zeta(x;\beta_1,\dots,\beta_{m-1},\alpha_i)\zeta(x;\alpha_1,\dots,\hat{\alpha}
_i,\dots,\alpha_{m+1})=0,\qquad  m\ge2,
\label{eq:5}
\end{gather}
where
\begin{gather*}
\zeta(x;\alpha_1,\dots,\alpha_{n})=\Delta(\alpha_1,\dots,\alpha_{n})\tau(x+[\alpha_1]+\dots+[\alpha_{n}]),
\\
\Delta(\alpha_1,\dots,\alpha_n)=\prod_{i<j}(\alpha_i-\alpha_j),
\end{gather*}
and $\hat{\alpha}_i$ means that $\alpha_i$ should be removed.
\begin{example}\label{example1}
In the case of $m=2$, we have
\begin{gather}
\alpha_{12}\alpha_{34}\tau(x+[\alpha_1]+[\alpha_2])\tau(x+[\alpha_3]+[\alpha_4])
-\alpha_{13}\alpha_{24}\tau(x+[\alpha_1]+[\alpha_3])\tau(x+[\alpha_2]+[\alpha_4])
\nonumber
\\
\qquad
{}+\alpha_{14}\alpha_{23}\tau(x+[\alpha_1]+[\alpha_4])\tau(x+[\alpha_2]+[\alpha_3])=0,
\label{eq:6}
\end{gather}
where $\alpha_{ij}=\alpha_i-\alpha_j$.
\end{example}

We call~\eqref{eq:6} `the three term equation'.
We have derived~\eqref{eq:6} from~\eqref{eq:1}.
The fact that the converse is true is proved by Takasaki and Takebe~\cite{TT1}.
\begin{theorem}[\cite{TT1}] \label{theorem1} The three term equation~\eqref{eq:6} is equivalent to the KP hierarchy~\eqref{eq:1}.
\end{theorem}

In~\cite{TT1} the theorem is proved by constructing the wave function of the KP-hierarchy.
To do it the dif\/ferential Fay identity, which is a~certain limit of~\eqref{eq:6}, is used.
Here we give an alternative and direct proof of the theorem.
\begin{proposition}\label{proposition1}
The KP hierarchy~\eqref{eq:1} is equivalent to~\eqref{eq:5}.
\end{proposition}
\begin{proof}
It is suf\/f\/icient to prove that if~\eqref{eq:4} holds for any $m\ge2$ and arbitrary
$\alpha_i\neq0$, $\beta_i$, then~\eqref{eq:3} holds.
Set the left hand side of~\eqref{eq:3} to be $F(y)$.
Expand $F(y)$ in $y$ as
\begin{gather*}%\label{eq:7}
F(y)=\sum_\gamma F_\gamma y^\gamma,\qquad  y^\gamma=y_1^{\gamma_1}y_2^{\gamma_2}
\cdots,\qquad \gamma=(\gamma_1,\gamma_2,\dots).
\end{gather*}
We consider the case $\beta_i=0$, $1\le i\le m-1$ in~\eqref{eq:22} and set
\begin{gather*}%\label{eq:8}
y'=\sum_{i=1}^{m+1}[\alpha_i].
\end{gather*}
We prove $F_\gamma=0$ for any $\gamma$ if $F\big({-}\frac{y'}{2}\big)=0$ for any $m\ge2$.
We consider $m$ f\/ixed.
Let us def\/ine the weight of $y_i$ to be $i$ and ${\rm wt}\,y^\gamma=\sum\limits_{i=1}^{\infty}i\gamma_i$.
Decompose $F$ according to weights as
\begin{gather*}
F=F^{(0)}+F^{(1)}+F^{(2)}+\cdots,
%\label{eq:9}
\qquad
F^{(i)}=\sum_{{\rm wt}\,y^\gamma=i}F_\gamma y^\gamma.
%\label{eq:10}
\end{gather*}
We substitute $-\frac{y'}{2}$ to $y$ and get the homogeneous polynomial of degree $i$ of
$\alpha_1,\dots,\alpha_{m+1}$:
\begin{gather*}%\label{eq:11}
F^{(i)}\left(-\frac{y'}{2}\right)=\sum_{\gamma_1+\dots+\gamma_{m+1}=i}b^{(i)}_{\gamma_1\dots\gamma_{m+1}}
\alpha_1^{\gamma_1}\cdots\alpha_{m+1}^{\gamma_{m+1}}.
\end{gather*}
Then $F\big({-}\frac{y'}{2}\big)=0$ is equivalent to $F^{(i)}\big({-}\frac{y'}{2}\big)=0$ for any $i$.
Notice that $i y'_i=\alpha_1^i+\cdots+\alpha_{m+1}^i$ is a~power sum symmetric function.
Therefore $y_1',\dots,y_{m+1}'$ are algebraically independent (see~\cite[(2.12)]{Mac}).
If $i\le m+1$, then $F^{(i)}(y)$ is a~polynomial at most of $y_1,\dots,y_{m+1}$.
Thus $F^{(i)}\big({-}\frac{y'}{2}\big)=0$ implies $F_\gamma=0$ for any $\gamma$ satisfying ${\rm wt}\,\gamma=i$.
Since $m$ is arbitrary we have $F_\gamma=0$ for any $\gamma$.
\end{proof}
\setcounter{section}{3}
\setcounter{subsection}{0}
\setcounter{subsubsection}{0}
\setcounter{Theorem}{0}
\setcounter{equation}{27}
\setcounter{theorem}{2}
\setcounter{proposition}{6}
\setcounter{example}{2}
% Original source lines 700--923
\section{BKP hierarchy}\label{section4}

Let $\tau(t)$ be the $\tau$-function of the BKP hierarchy.
In this case, the time variable is $t=(t_1,t_3,t_5,\dots)$.
We set
\begin{gather*}
[\alpha]_{\rm o}=\left(\alpha,\frac{\alpha^3}{3},\frac{\alpha^5}{5},\dots\right),\hskip1cm\tilde{\xi}
(t,\lambda)=\sum_{n=1}^{\infty}t_{2n-1}\lambda^{2n-1}.
\end{gather*}
The BKP hierarchy \cite{DJKM1,JM1} is def\/ined by
\begin{gather}
\oint e^{\tilde{\xi}(t-t',\lambda)}\tau\big(t-2\big[\lambda^{-1}\big]_{\rm o}\big)\tau\big(t'+2\big[\lambda^{-1}\big]_{\rm o}\big)\frac{d\lambda}
{2\pi i\lambda}=\tau(t)\tau(t').
\label{eq:50}
\end{gather}
Set $t=x-y$, $t'=x+y$.
We get
\begin{gather}
\oint e^{-2\tilde{\xi}(y,\lambda)}\tau\big(x-y-2\big[\lambda^{-1}\big]_{\rm o}\big)\tau\big(x+y+2\big[\lambda^{-1}\big]_{\rm o}\big)\frac{d\lambda}
{2\pi i\lambda}=\tau(x-y)\tau(x+y).
\label{eq:51}
\end{gather}
Set
\begin{gather*}
y=\sum_{i=1}^{n} [\alpha_i]_{\rm o}
\end{gather*}
in~\eqref{eq:51}. Then we have
\begin{gather*}%\label{eq:52}
\oint e^{-2\tilde{\xi}\big(\sum\limits_{i=1}^{n} [\alpha_i]_{\rm o},\lambda\big)}\tau\left(x-\sum_{i=1}^{n}
 [\alpha_i]_{\rm o}-2\big[\lambda^{-1}\big]_{\rm o}\right)
\tau\left(x+\sum_{i=1}^{n} [\alpha_i]_{\rm o}+2\big[\lambda^{-1}\big]_{\rm o}\right)\frac{d\lambda}{2\pi i\lambda}
\\
\qquad
{}=\tau\left(x-\sum_{i=1}^{n} [\alpha_i]_{\rm o}\right)\tau\left(x+\sum_{i=1}^{n} [\alpha_i]_{\rm o}\right).
\end{gather*}
By decomposing $-2\sum\limits_{n=1}^\infty t_{2n-1}{\lambda}^{2n-1}$ as
\begin{gather*}
-2\sum_{n=1}^\infty t_{2n-1}{\lambda}^{2n-1}=-\sum_{n=1}^{\infty}t_n{\lambda}^n+\sum_{n=1}^{\infty}
t_n(-\lambda)^n,
\end{gather*}
we get
\begin{gather*}
\exp\left(-2\tilde{\xi}\left(\sum_{i=1}^{n} [\alpha_i]_{\rm o},\lambda\right)\right)=\prod_{i=1}^{n}
\frac{1-\alpha_i\lambda}{1+\alpha_i\lambda}.
\end{gather*}
Computing the integral by taking residues as before, shifting $x$ as $x+\sum\limits_{i=1}^{n}[\alpha_i]_{\rm o}$ and
divi\-ding by~$\tau(x)^2$ we have
\begin{gather}
\sum_{i=1}^{n}(-1)^{i-1}\frac{\tau(x+2[\alpha_i]_{\rm o})}{\tau(x)}A_{1\dots\hat{i}\dots n}^{-1}
\frac{\tau\left(x+2\sum\limits_{l=1,\, l\neq i}^n[\alpha_l]_{\rm o}\right)}{\tau(x)}\nonumber
\\
\qquad
{}-A_{1\dots n}^{-1}\frac{\tau\left(x+2\sum\limits_{l=1}^{n}[\alpha_l]_{\rm o}\right)}{\tau(x)}=0,\qquad  n \ \text{odd},
\label{eq:53}
\\
\sum_{i=1}^{n-1}(-1)^{i-1}\frac{\alpha_{i,n}}{\tilde{\alpha}_{i,n}}
\frac{\tau(x+2[\alpha_i]_{\rm o}+2[\alpha_n]_{\rm o})}{\tau(x)}A_{1\dots\hat{i}\dots n-1}^{-1}
\frac{\tau\left(x+2\sum\limits_{l=1,\, l\neq i}^{n-1}[\alpha_l]_{\rm o}\right)}{\tau(x)}\nonumber
\\
\qquad
{} -A_{1\dots n}^{-1}\frac{\tau\left(x+2\sum\limits_{l=1}^{n}[\alpha_l]_{\rm o}\right)}{\tau(x)}=0,\qquad  n \ \text{even}.
\label{eq:54}
\end{gather}
Here $A_{1\dots n}$ is def\/ined by
\begin{gather*}
A_{1\dots n}=\prod_{i<j}^{n}\frac{\tilde{\alpha}_{ij}}{\alpha_{ij}},\qquad \tilde{\alpha}_{ij}
=\alpha_i+\alpha_j,\qquad \alpha_{ij}=\alpha_i-\alpha_j.
\end{gather*}
\begin{example}\label{example3}
The case $n=3$ of~\eqref{eq:53} becomes
\begin{gather}
\frac{\tau\left(x+2\sum\limits_{i=1}^{3}[\alpha_i]_{\rm o}\right)}{\tau(x)}=A_{123}\left\{\frac{\tau(x+2[\alpha_1]_{\rm o})}{\tau(x)}
\frac{\alpha_{23}}{\tilde{\alpha}_{23}}\frac{\tau(x+2[\alpha_2]_{\rm o}+2[\alpha_3]_{\rm o})}{\tau(x)}\right.
\nonumber
\\
\hskip36mm-\frac{\tau(x+2[\alpha_2]_{\rm o})}{\tau(x)}\frac{\alpha_{13}}{\tilde{\alpha}_{13}}
\frac{\tau(x+2[\alpha_1]_{\rm o}+2[\alpha_3]_{\rm o})}{\tau(x)}\nonumber
\\
\hskip36mm\left.
+\frac{\tau(x+2[\alpha_3]_{\rm o})}{\tau(x)}\frac{\alpha_{12}}{\tilde{\alpha}_{12}}
\frac{\tau(x+2[\alpha_1]_{\rm o}+2[\alpha_2]_{\rm o})}{\tau(x)}\right\}.
\label{eq:55}
\end{gather}
\end{example}
We call equation~\eqref{eq:55} `the four term equation of the BKP hierarchy'.
\begin{example}
The case of $n=4$ of~\eqref{eq:54} is
\begin{gather}
\frac{\tau\left(x+2\sum\limits_{i=1}^{4}[\alpha_i]_{\rm o}\right)}{\tau(x)}
=A_{1234}\left\{\frac{\alpha_{14}}{\tilde{\alpha}_{14}}
\frac{\tau(x+2[\alpha_1]_{\rm o}+2[\alpha_4]_{\rm o})}{\tau(x)}\frac{\alpha_{23}}{\tilde{\alpha}_{23}}
\frac{\tau(x+2[\alpha_2]_{\rm o}+2[\alpha_3]_{\rm o})}{\tau(x)}\right.
\nonumber
\\
\hphantom{\frac{\tau\left(x+2\sum\limits_{i=1}^{4}[\alpha_i]_{\rm o}\right)}{\tau(x)}=}{}
-\frac{\alpha_{24}}{\tilde{\alpha}_{24}}\frac{\tau(x+2[\alpha_2]_{\rm o}+2[\alpha_4]_{\rm o})}{\tau(x)}
\frac{\alpha_{13}}{\tilde{\alpha}_{13}}\frac{\tau(x+2[\alpha_1]_{\rm o}+2[\alpha_3]_{\rm o})}{\tau(x)}\nonumber
\\
\left.
\hphantom{\frac{\tau\left(x+2\sum\limits_{i=1}^{4}[\alpha_i]_{\rm o}\right)}{\tau(x)}=}{}
+\frac{\alpha_{34}}{\tilde{\alpha}_{34}}\frac{\tau(x+2[\alpha_3]_{\rm o}+2[\alpha_4]_{\rm o})}{\tau(x)}\frac{\alpha_{12}
}{\tilde{\alpha}_{12}}\frac{\tau(x+2[\alpha_1]_{\rm o}+2[\alpha_2]_{\rm o})}{\tau(x)}\right\}.
\label{eq:56}
\end{gather}
\end{example}

 As is proved in Proposition~\ref{proposition8}, equation~\eqref{eq:56} can be derived from equation~\eqref{eq:55}.
\begin{theorem}\label{theorem3}
The four term equation~\eqref{eq:55} is equivalent to the BKP hierarchy~\eqref{eq:50}.
\end{theorem}

We prove this theorem in a~similar way to the case of the KP hierarchy.

In order to prove the theorem, we use Pfaf\/f\/ians.
The def\/inition and notation of Pfaf\/f\/ians are reviewed in Appendix~\ref{appendixA}.

Let us def\/ine the components of Pfaf\/f\/ians by
\begin{gather*}%\label{eq:57}
(0,j)=\frac{\tau(x+2[\alpha_j]_{\rm o})}{\tau(x)},\qquad \hskip1cm(i,j)=\frac{\alpha_{ij}}{\tilde{\alpha}_{ij}}
\frac{\tau(x+2[\alpha_i]_{\rm o}+2[\alpha_j]_{\rm o})}{\tau(x)}.
\end{gather*}
Then it is possible to rewrite~\eqref{eq:55} and~\eqref{eq:56} as
\begin{gather}
\frac{\tau\left(x+2\sum\limits_{i=1}^{3}[\alpha_i]_{\rm o}\right)}{\tau(x)}=A_{123}(0,1,2,3),
\label{eq:58}
\\
\frac{\tau\left(x+2\sum\limits_{i=1}^{4}[\alpha_i]_{\rm o}\right)}{\tau(x)}=A_{1234}(1,2,3,4),
\label{eq:59}
\end{gather}
respectively.

The following proposition can be proved in a~similar manner to Proposition~\ref{proposition1}.
\begin{proposition}\label{proposition7}
The BKP hierarchy~\eqref{eq:50} is equivalent to~\eqref{eq:53} and~\eqref{eq:54}.
\end{proposition}
\begin{proposition}\label{proposition8}
The following equations are implied by~\eqref{eq:55}:
\begin{gather}
\frac{\tau\left(x+2\sum\limits_{i=1}^{n}[\alpha_i]_{\rm o}\right)}{\tau(x)}=A_{1\dots n}(0,1,2,\dots,n),\qquad n\ge3, \ \text{odd},
\label{eq:66}
\\
\frac{\tau\left(x+2\sum\limits_{i=1}^{n}[\alpha_i]_{\rm o}\right)}{\tau(x)}=A_{1\dots n}(1,2,\dots,n),\qquad n\ge4, \ \text{even}.
\label{eq:67}
\end{gather}
\end{proposition}

\begin{proof}
First we prove that~\eqref{eq:55} implies~\eqref{eq:56}.
Shift $x$ in~\eqref{eq:55} as $x\rightarrow x+2[\alpha_4]_{\rm o}$:
\begin{gather}
\frac{\tau\left(x+2\sum\limits_{i=1}^{4}[\alpha_i]_{\rm o}\right)}{\tau(x+2[\alpha_4]_{\rm o})}\nonumber\\
\qquad{}
=A_{123}\frac{\tau(x)^2}{\tau(x+2[\alpha_4]_{\rm o})^2}\left\{\frac{\tau(x+2[\alpha_1]_{\rm o}+2[\alpha_4]_{\rm o})}{\tau(x)}
\frac{\alpha_{23}}{\tilde{\alpha}_{23}}\frac{\tau(x+2[\alpha_2]_{\rm o}+2[\alpha_3]_{\rm o}+2[\alpha_4]_{\rm o})}{\tau(x)}
\right.
\nonumber
\\
\qquad\quad{} -\frac{\tau(x+2[\alpha_2]_{\rm o}+2[\alpha_4]_{\rm o})}{\tau(x)}\frac{\alpha_{13}}{\tilde{\alpha}_{13}}
\frac{\tau(x+2[\alpha_1]_{\rm o}+2[\alpha_3]_{\rm o}+2[\alpha_4]_{\rm o})}{\tau(x)}\nonumber
\\
 \left.
 \qquad\quad{} +\frac{\tau(x+2[\alpha_3]_{\rm o}+2[\alpha_4]_{\rm o})}{\tau(x)}\frac{\alpha_{12}}{\tilde{\alpha}_{12}}
\frac{\tau(x+2[\alpha_1]_{\rm o}+2[\alpha_2]_{\rm o}+2[\alpha_4]_{\rm o})}{\tau(x)}\right\}.
\label{eq:68}
\end{gather}
Use~\eqref{eq:58} to rewrite $\tau(x+2[\alpha_{i_1}]_{\rm o}+2[\alpha_{i_2}]_{\rm o}+2[\alpha_{i_3}]_{\rm o})$
in~\eqref{eq:68}.
Then we get~\eqref{eq:59}.
We prove~\eqref{eq:66} by induction on $n$.
The case $n=3$ is obvious.
Suppose that~\eqref{eq:66} holds in case of $n$:
\begin{gather}
\frac{\tau\left(x+2\sum\limits_{i=1}^{n}[\alpha_i]_{\rm o}\right)}{\tau(x)}=A_{1\dots n}(0,1,2,\dots,n)=A_{1\dots n}{\rm Pf}\, {\mathcal A},
\label{eq:69}
\end{gather}
where ${\mathcal A}=(a_{ij})_{0\le i,j\le n}$ is a~skew-symmetric matrix,
\begin{gather*}%\label{eq:70}
a_{ij}=
\begin{cases}\dfrac{\tau(x+2[\alpha_j]_{\rm o})}{\tau(x)}, &  i=0,
\vspace{2mm}\\
\dfrac{\alpha_{ij}}{{\tilde{\alpha}}_{ij}}\dfrac{\tau(x+2[\alpha_i]_{\rm o}+2[\alpha_j]_{\rm o})}{\tau(x)}
,&  i\neq0,\quad  i<j.
\end{cases}
\end{gather*}
In~\eqref{eq:69} we shift $x$ as
\begin{gather*}
x\rightarrow x+2[\alpha_{n+1}]_{\rm o}+2[\alpha_{n+2}]_{\rm o}.
\end{gather*}
Then we have, using~\eqref{eq:58} and~\eqref{eq:59},
\begin{gather}
\frac{\tau\left(x+2\sum\limits_{i=1}^{n+2}[\alpha_i]_{\rm o}\right)}{\tau(x)}=A_{1\dots n+2}(n+1,n+2)^{-\frac{n-1}{2}}{\rm Pf}\,
{\mathcal B},
\label{eq:71}
\end{gather}
where ${\mathcal B}=(b_{ij})_{0\le i<j\le n}$ is a~skew-symmetric matrix given by
\begin{gather*}%\label{eq:72}
b_{ij}=(n+1,n+2,i,j),\qquad  i<j .
\end{gather*}

By the analogue of the Sylvester' theorem for Pfaf\/f\/ians (Appendix~\ref{appendixB}), we have
\begin{gather*}%\label{eq:73}
{\rm Pf}((1,2,\dots,2r,i,j))_{2r+1\le i<j\le2m}=(1,2,\dots,2r)^{m-r-1}(1,2,\dots,2m).
\end{gather*}
Consider the case $r=1$ and $2m=n+3$ of this formula:
\begin{gather}
{\rm Pf}((n+1,n+2,i,j))_{0\le i<j\le n}=(n+1,n+2)^{\frac{n-1}{2}}(n+1,n+2,0,\dots,n).
\label{eq:80}
\end{gather}
Substituting~\eqref{eq:80} to~\eqref{eq:71}, we get
\begin{gather*}%\label{eq:74}
\frac{\tau\left(x+2\sum\limits_{i=1}^{n+2}[\alpha_i]_{\rm o}\right)}{\tau(x)}=A_{1\dots n+2}(0,1,\dots,n+2).
\end{gather*}
The case of even $n$ is similarly proved.
\end{proof}

\begin{proposition}\label{proposition9}
The Pl\"{u}cker's relations for Pfaffians of the right hand side of~\eqref{eq:66} and~\eqref{eq:67}
give the addition formulae~\eqref{eq:53} and~\eqref{eq:54} respectively.
\end{proposition}

\begin{proof}
Using the Pl\"{u}cker's relations~\eqref{eq:500} and~\eqref{eq:501} for Pfaf\/f\/ians given in Appendix~\ref{appendixC},
the proposition can easily be checked by direct calculations.
\end{proof}
\setcounter{section}{4}
\setcounter{subsection}{0}
\setcounter{subsubsection}{0}
\setcounter{Theorem}{0}
\setcounter{equation}{41}
% Original source lines 925--1022
\appendix

\section{Pfaf\/f\/ians} \label{appendixA}

Let ${\mathcal A} =(a_{ij})_{1\le i,j\le2m}$ be a~skew-symmetric matrix.
Then the Pfaf\/f\/ian ${\rm Pf}\, {\mathcal A}$ is def\/ined by
\begin{gather*}
\det{\mathcal A}=({\rm Pf}\,{\mathcal A})^2,\qquad {\rm Pf}\,{\mathcal A}=a_{12}a_{34}\cdots a_{2m-1,2m}+\cdots.
\end{gather*}
Following~\cite{H1} we denote ${\rm Pf}\, {\mathcal A}$ by $(1,2,3,\dots,2m)$:
\begin{gather*}
{\rm Pf}\,{\mathcal A}=(1,2,3,\dots,2m).
\end{gather*}
It is directly def\/ined by
\begin{gather*}
(1,2,3,\dots,2m)=\sum{\rm sgn}(i_1,\dots,i_{2m})\cdot(i_1,i_2)(i_3,i_4)\cdots(i_{2m-1},i_{2m}
),\qquad (i,j)=a_{ij},
\end{gather*}
where the sum is over all permutations of $(1,\dots,2m)$ such that
\begin{gather*}
i_1<i_3<\cdots<i_{2m-1},\qquad i_1<i_2,\dots,i_{2m-1}<i_{2m},
\end{gather*}
and ${\rm sgn}(i_1,\dots,i_{2m})$ is the signature of the permutation $(i_1,\dots,i_{2m})$.
The Pfaf\/f\/ian can be expanded as
\begin{gather*}
(1,2,3,\dots,2m)=\sum_{j=2}^{2m}(-1)^j(1,j)(2,3,\dots,\hat{j},\dots,2m).
\end{gather*}
For example, in the case of $m=2$,
\begin{gather*}
(1,2,3,4)=(1,2)(3,4)-(1,3)(2,4)+(1,4)(2,3).
\end{gather*}

\section{Sylvester's theorem for determinants and Pfaf\/f\/ians} \label{appendixB}

The following theorem is known as
Sylvester's theorem.
\begin{theorem}\label{theorem4}
Let $r\le m$, ${\mathcal A}=(a_{ij})_{1\le i,j\le m}$ and ${\mathcal A}_r=(a_{ij})_{1\le i,j\le r}$.
Set
\begin{gather*}
{\mathcal B}=(b_{ij})_{r+1\le i,j\le m},
\qquad
b_{ij}=\det\left(
\begin{matrix}a_{11}&\hdots&a_{1r}&a_{1j}
\\
\vdots&\ddots&\vdots&\vdots
\\
a_{r1}&\hdots&a_{rr}&a_{rj}
\\
a_{i1}&\hdots&a_{ir}&a_{ij}
\end{matrix}
\right).
\end{gather*}
Then we get
\begin{gather*}%\label{eq:90}
\det{\mathcal B}=(\det{\mathcal A}_r)^{m-r-1}\det{\mathcal A}.
\end{gather*}
\end{theorem}

  Let ${\mathcal A}=(a_{ij})_{1\le i,j\le2m}$ be a~skew-symmetric matrix and set $(i,j)=a_{ij}$.
For $r\le m$ let ${\mathcal P} =(p_{ij})_{2r+1\le i,j\le2m}$, $p_{ij}=(1,2,\dots,2r,i,j)$ and $I_r=\{1,2,\dots,2r\}$.
In general, for a~subset $I\subset\{1,2,\dots,2m\}$ we set ${\mathcal A}(I)=(a_{ij})_{i,j\in I}$ and for $i<j$, $k<l$
we denote by ${\mathcal A}_{kl}^{ij}$ be the square matrix of degree $2(m-1)$ which is obtained from ${\mathcal A}$ by
removing $i$-th and $j$-th rows, $k$-th and $l$-th columns.

\begin{theorem}[\cite{H2}]\label{theorem5} For $r\le m$ the following identity holds:
\begin{gather*}%\label{eq:91}
{\rm Pf}\,{\mathcal P}=({\rm Pf}\,{\mathcal A}(I_r))^{m-r-1}{\rm Pf}\,{\mathcal A}.
\end{gather*}
\end{theorem}

\section{The Pl\"{u}cker relation for Pfaf\/f\/ians}\label{appendixC}

 There exist analogues of the Pl\"{u}cker's
relations for Pfaf\/f\/ians~\cite{O1}.
They are given by
\begin{gather}
\sum_{l=1}^L(-1)^l(i_1,\dots,i_K,j_l)(j_1,\dots,\hat{j}_l,\dots,j_L)
\nonumber
\\
\qquad
{}+\sum_{k=1}
^K(-1)^k(i_1,\dots,\hat{i}_k,\dots,i_K)(j_1,\dots,j_L,i_k)=0,
\label{eq:75}
\end{gather}
where $K$ and $L$ are odd.
We understand that $(\varnothing)=1$.

For $n$ odd, taking $K=1$, $L=n$, $i_1=0$ and $j_1,\dots,j_n\neq0$ in~\eqref{eq:75}, we get
\begin{gather}
\sum_{l=1}^n(-1)^{l-1}(0,j_l)(j_1,\dots,\hat{j}_l,\dots,j_n)-(0,j_1,\dots,j_n)=0.
\label{eq:500}
\end{gather}
For $n$ even, setting $K=1$, $L=n-1$ and $i_1\neq0$ in~\eqref{eq:75}, we have
\begin{gather}
\sum_{l=1}^{n-1}(-1)^{l-1}(i_1,j_l)(j_1,\dots,\hat{j}_l,\dots,j_{n-1})-(i_1,j_1,\dots,j_{n-1})=0.
\label{eq:501}
\end{gather}
\begin{thebibliography}{99}
\bibitem[6]{DJKM1}
Date E., Kashiwara M., Jimbo M., Miwa T., Transformation groups for soliton
  equations, in Nonlinear In\-tegrable Systems~-- Classical Theory and Quantum
  Theory ({K}yoto, 1981), World Sci. Publishing, Singapore, 1983, 39--119.

\bibitem[9]{H2}
Hirota R., Generalizations of determinant identities by Pfaf\/f\/ian, in
  Mathematical Theories and Applications of Nonlinear Waves and Nonlinear
  Dynamics, Research Institute for Applied Mechanics, Kyushu University, 2004, 148--156.

\bibitem[10]{H1}
Hirota R., The direct method in soliton theory, \href{http://dx.doi.org/10.1017/CBO9780511543043}{\textit{Cambridge Tracts in
  Mathematics}}, Vol.~155, Cambridge University Press, Cambridge, 2004.


\bibitem[12]{JM1}
Jimbo M., Miwa T., Solitons and inf\/inite-dimensional {L}ie algebras,
  \href{http://dx.doi.org/10.2977/prims/1195182017}{\textit{Publ. Res. Inst. Math. Sci.}} \textbf{19} (1983), 943--1001.

\bibitem[13]{MJD1}
Miwa T., Jimbo M., Date E.,
Solitons. Dif\/ferential equations, symmetries and inf\/inite-dimensional algebras,
\textit{Cambridge Tracts in Mathematics}, Vol.~135, Cambridge University Press, Cambridge, 2000.

\bibitem[14]{Mac}
Macdonald I.G., Symmetric functions and {H}all polynomials, 2nd ed., \textit{Oxford
  Mathematical Monographs}, The Clarendon Press, Oxford University Press, New
  York, 1995.

\bibitem[15]{M1}
Miwa T., On {H}irota's dif\/ference equations, \href{http://dx.doi.org/10.3792/pjaa.58.9}{\textit{Proc. Japan Acad. Ser.~A
  Math. Sci.}} \textbf{58} (1982), 9--12.

\bibitem[18]{O1}
Ohta Y., Bilinear theory of solitons with {P}faf\/f\/ian labels,
  \textit{S\=urikaisekikenky\=usho K\=oky\=uroku}  (1993), no.~822, 197--205.

\bibitem[20]{SS1}
Sato M., Sato Y., Soliton equations as dynamical systems on inf\/inite
  dimensional Grassmann manifold, in Nonlinear PDE in Applied Science,
  \textit{North-Holland Math. Stud.}, Vol.~81, Editors H.~Fujita, P.~Lax,
  G.~Strang, Tokyo, 1982, 259--271.

\bibitem[23]{TT1}
Takasaki K., Takebe T., Integrable hierarchies and dispersionless limit,
  \href{http://dx.doi.org/10.1142/S0129055X9500030X}{\textit{Rev. Math. Phys.}} \textbf{7} (1995), 743--808,
  \href{http://arxiv.org/abs/hep-th/9405096}{hep-th/9405096}.

\end{thebibliography}
\end{document}
